
\section{Phase-polynomial states and circuits}
\label{sec:phase-polynomial-states}

Our simulation framework rests on a simple description of states prepared by Pauli, $\CNOT$, and $\T$ gates: a uniform superposition over an affine space, with phases specified by a third-order polynomial.
We first introduce the CNOT-dihedral group and its phase-polynomial description, then define the corresponding class of third-order phase-polynomial (3PP) states.
Finally, we introduce compatible measurements, which preserve this state class, and use them to specify the circuits considered throughout the paper.

\subsection{The CNOT-dihedral group and phase polynomials}

We consider circuits generated by Pauli, $\CNOT$, and $\T$ gates, whose unitary operations form the third-order CNOT-dihedral group \cite{cross2016scalable}.

\begin{definition}[Third-order CNOT-dihedral group]
    The third-order CNOT-dihedral group on $n$ qubits is
    \begin{align}
        G_3 = \langle X_j, \CNOT_{ij}, \T_j : i \neq j \rangle / \langle \omega_8 I \rangle,
    \end{align}
    where $\omega_8 = e^{i\pi/4}$ and $\T = \operatorname{diag}(1,\omega_8)$.
    The quotient identifies operators that differ by a power of the global phase $\omega_8$.
\end{definition}

Throughout this work, we refer to $G_3$ simply as the \emph{CNOT-dihedral group}.
The name comes from the fact that $X$ and $\T$ generate a dihedral group up to global phase, to which $\CNOT$ gates are added.
Since $Z=\T^4$ and $Y=iXZ$, this group contains all Pauli operators up to global phase.

CNOT-dihedral operators admit an efficient description through their action on computational-basis states.
The $X$ and $\CNOT$ gates implement invertible affine transformations of the bit strings, while $\T$ gates contribute phases that are powers of $\omega_8$.
Consequently, every $U\in G_3$ has a unitary representative satisfying
\begin{align}
    U\ket{x} = \omega_8^{p(x)}\ket{Ax\oplus b},
    \qquad x\in\mathbb{F}_2^n,
\end{align}
where $A\in\mathrm{GL}(n,\mathbb{F}_2)$, $b\in\mathbb{F}_2^n$, and $p:\mathbb{F}_2^n\to\mathbb{Z}_8$.
The phase functions $p$ arising in this way are precisely the third-order phase polynomials defined below \cite{amy2013meetinthemiddle,campbell2017unified,amy2019tcount,amy2018finite}.

\begin{definition}[Third-order phase polynomial]\label{def:third-order-phase-polynomial}
A function $p:\F_2^n\to\Z_8$ is a \emph{third-order phase polynomial} if it can be written as
\begin{align}
    p(x) ={}& d + \sum_{i=1}^n a_i x_i
        + 2 \sum_{1 \le i < j \le n} b_{ij} x_i x_j \notag \\
    &+ 4 \sum_{1 \le i < j < k \le n} c_{ijk} x_i x_j x_k,
\end{align}
with $d,a_i \in \Z_8$, $b_{ij} \in \Z_4$, and $c_{ijk} \in \Z_2$, and where arithmetic is modulo $8$.
\end{definition}

This expression is sometimes called a \emph{weighted polynomial}~\cite{campbell2017unified,vuillot2022quantum}, while the term ``phase polynomial'' is often reserved for expressions of the form
\begin{align} \label{eqn:phase-polynomial-def}
    p(x)=a_0+\sum_{u\ne0}\lambda_u\langle u,x\rangle\pmod8,
\end{align}
where the inner product is taken modulo $2$, and the rest is taken modulo $8$. We show in \Cref{thm:phase-polynomial-equivalence} that these two forms are equivalent, and we therefore use the terms interchangeably.

Within this class, diagonal Clifford operators correspond to a more restricted form.

\begin{definition}[Clifford phase polynomial]\label{def:clifford-phase-polynomial}
A function $d:\F_2^n\to\Z_8$ is a \emph{Clifford phase polynomial} if it has the form
\begin{align}
    d(x)=d_0+2\sum_i \alpha_i x_i
        +4\sum_{i<j}\beta_{ij}x_ix_j \pmod 8,
\end{align}
where $d_0\in\Z_8$, $\alpha_i\in\Z_4$, and $\beta_{ij}\in\F_2$.
\end{definition}

These are precisely the phase functions of diagonal Clifford operators, up to global phase. \Cref{thm:clifford-phase-polynomial} proves this characterization and gives the equivalent parity-sum form.

\subsection{Third-order phase-polynomial states}

The efficient representation of CNOT-dihedral operators extends to the states they prepare. Indeed, for $U\in G_3$ on $n=m+\ell$ qubits,
\begin{align}
    U\ket{+}^{\otimes m} \ket{0}^{\otimes \ell} = \frac{1}{\sqrt{2^m}}\sum_{x \in \F_2^m} \omega_8^{p((x,0))} \ket{A(x,0) \oplus b}
\end{align}
where $(x,0)$ is the concatenation of $x \in \F_2^m$ and $0 \in \F_2^\ell$. Writing $\tilde{p}(x):= p((x,0))$ and $\tilde{A}$ for the first $m$ columns of $A$, we can rewrite the above equation as
\begin{align}
    U\ket{+}^{\otimes m} \ket{0}^{\otimes \ell} = \frac{1}{\sqrt{2^m}}\sum_{x \in \F_2^m} \omega_8^{\tilde{p}(x)} \ket{\tilde{A}x \oplus b}.
\end{align}
Since $\tilde{A}$ has full column rank, choose a left inverse $L$ satisfying $L\tilde{A}=I_m$ and define $q(y):=\tilde p(L(y\oplus b))$ for every $y\in\F_2^n$. Then $q(\tilde{A}x\oplus b)=\tilde p(x)$, and $q$ is a third-order phase polynomial by \Cref{cor:phase-polynomial-affine-pullback}. Changing variables therefore gives
\begin{align}
    U\ket{+}^{\otimes m} \ket{0}^{\otimes \ell} = \frac{1}{\sqrt{|S|}}\sum_{x \in S} \omega_8^{q(x)} \ket{x},
\end{align}
where $S = b \oplus \Im(\tilde{A})$. This motivates the following state class.

\begin{definition}[3PP states]\label{def:3pp}
    A \emph{third-order phase-polynomial state} (3PP state) is a state that is, up to global phase, of the form
    \begin{align}
        \ket{\psi} = \frac{1}{\sqrt{|S|}} \sum_{x \in S} \omega_8^{q(x)} \ket{x},
    \end{align}
    where $S \subseteq \F_2^n$ is a nonempty affine subspace and $q:\F_2^n\to\Z_8$ is a third-order phase polynomial.
\end{definition}

We note that the polynomial associated to a 3PP state is not unique, since its value only matters on the affine support $S$. We present later in \cref{thm:efficient-simulation} an algorithm to efficiently determine whether two 3PP states, represented by their phase polynomial and support, are equal up to a global phase.

The preparation argument above also has a converse.

\begin{proposition}[Preparation of 3PP states]\label{prop:3pp-preparation}
    A state is a 3PP state if and only if, up to global phase, it can be prepared from $\ket{0}$ and $\ket{+}$ states by a circuit of Pauli, $\CNOT$, and $\T$ gates.
\end{proposition}
\begin{proof}
    The calculation above proves that every state prepared by such a circuit is a 3PP state. Conversely, let $\ket{\psi}$ have the form in \Cref{def:3pp}. A uniform superposition over the affine space $S$ can be prepared from $\ket{0}$ and $\ket{+}$ states using $X$ and $\CNOT$ gates: choose an affine parametrization $S=b\oplus\Im(B)$ with $B$ of full column rank, extend the columns of $B$ to an invertible binary matrix, and implement the resulting affine transformation. Writing $q$ in the weighted form above and applying
    \[
        \prod_i \T_i^{a_i}
        \prod_{i<j}\CS_{ij}^{b_{ij}}
        \prod_{i<j<k}\CCZ_{ijk}^{c_{ijk}}
    \]
    then imprints the phase $\omega_8^{q(x)}$ up to the global phase $\omega_8^d$.
    Since $\CS$ and $\CCZ$ can themselves be synthesized from $\T$ and $\CNOT$ gates, this gives the required preparation.
\end{proof}

Keeping track of the affine space $S$ and the $O(n^3)$ coefficients of $q$ is therefore enough to represent any 3PP state.

\subsection{Compatible measurements and 3PP circuits}\label{sec:compatible-measurements}

To extend the gate set to measurements, we require each measurement branch to remain within the 3PP state class.

\begin{definition}[Compatible measurement]\label{def:compatible-measurement}
    Given a 3PP state $\ket{\psi}$, we say that a projective measurement of a Pauli operator is \emph{compatible} with it if the post-measurement state is a 3PP state for every outcome with nonzero probability.
\end{definition}

The need for introducing the notion of compatibility comes from the observation that not every $X$ measurement is compatible, as shown in the following simple example.
\begin{example}\label{ex:incompatible-x-measurement}
    Consider the 3PP state $\ket{\CCZ}=\CCZ\ket{+}^{\otimes 3}$. For an outcome $\varepsilon\in\{+1,-1\}$ of measuring the first qubit in the $X$ basis, the unnormalized post-measurement vector is
    \begin{align}
        \ket{\widetilde\psi_\varepsilon}=\frac{1}{2}(1+\varepsilon X_1)\ket{\CCZ}.
    \end{align}
    The outcome $\varepsilon=-1$ leaves the 3PP state $\ket{\psi_-}\propto\ket{-}\otimes\ket{11}$. The outcome $\varepsilon=+1$, which also has nonzero probability, leaves
    \begin{align}
        \ket{\psi_+}\propto\ket{+}\otimes
        \bigl(\ket{00}+\ket{01}+\ket{10}\bigr).
    \end{align}
    Its computational-basis support has six elements, whereas every affine subspace over $\F_2$ has cardinality a power of two. Thus its support is not affine, so $\ket{\psi_+}$ is not a 3PP state and the measurement is incompatible.
\end{example}

\begin{figure}[tbp]
    \centering
    \begin{quantikz}
        \lstick{$\ket{\psi}$} & \ctrl{1} & \meter{X} & \cwbend{1} \\
        \lstick{$\ket{+}$}    & \control{} & \qw & \gate{X} & \qw
    \end{quantikz}
    \caption{Hadamard gadget. A $\CZ$ gate with a $\ket{+}$ ancilla, followed by an $X$ measurement of the input qubit and an outcome-dependent $X$ correction, leaves the ancilla in the state $\Hgate\ket{\psi}$.}
    \label{fig:hadamard-gadget}
\end{figure}

The fact that our framework cannot simulate arbitrary $X$ measurements is not surprising.
Indeed, preparing a $|+\rangle$ ancilla state, applying $\CZ$, measuring the input qubit in the $X$ basis, and applying an outcome-dependent $X$ correction to the ancilla, implements a Hadamard gate (see \cref{fig:hadamard-gadget}).
Allowing arbitrary $X$ measurements therefore gives universal quantum computation. Consequently, a polynomial-time classical sampler for these circuits would imply $\mathrm{BQP}\subseteq\mathrm{BPP}$.

By contrast, as we discuss in \cref{sec:z-measurements-3pp}, a $Z$-type Pauli measurement automatically preserves the 3PP form: conditioning on a nonzero-probability outcome intersects the affine support with an affine hyperplane and leaves the phase polynomial unchanged.

Together, these closure properties specify the circuits considered in our simulation algorithms.

\begin{definition}[3PP circuit]\label{def:3pp-circuit}
A \emph{3PP circuit} is composed of $\ket{0}$ and $\ket{+}$ preparations, Pauli, $\CNOT$, and $\T$ gates, arbitrary $Z$-type Pauli measurements, and $X$-type Pauli measurements that are compatible on every state reachable with nonzero probability. The operations may depend on earlier measurement outcomes through classical feedforward.
\end{definition}

Thus every measurement branch of a 3PP circuit can be described using an affine support and a third-order phase polynomial.
